% Part:sets-functions-relations
% Chapter: sets
% Section: reduction
% Hindi translation (hi-Deva-IN) of Open Logic commit 9620cc73f9c8e0ad003c514a5d3748f29611c4c0.

\documentclass[../../../include/open-logic-section]{subfiles}

\begin{document}

\olfileid[hi]{sfr}{siz}{red}

\olsection{न्यूनीकरण}

\begin{editorial}
  यह खंड न्यूनीकरण द्वारा अगणनीयता सिद्ध करता है और इसके परिणाम
  \olref[nen]{sec} के परिणामों से मेल खाते हैं। \olref[nen-alt]{sec} के
  परिणामों से मेल खाने वाला थोड़ा अधिक संक्षिप्त वैकल्पिक रूप
  \olref[red-alt]{sec} में दिया गया है।
\end{editorial}

हमने विकर्णीकरण के तर्क से दिखाया कि $\Pow{\PosInt}$ !!{nonenumerable} है।
हमारे पास पहले से यह प्रमाण था कि $\Bin^\omega$, अर्थात $0$ और $1$ के सभी
अनंत अनुक्रमों का समुच्चय, !!{nonenumerable} है। अब $\Pow{\PosInt}$ के
!!{nonenumerable} होने को सिद्ध करने का एक और ढंग देखें: दिखाइए कि \emph{यदि
$\Pow{\PosInt}$ !!{enumerable} है, तो $\Bin^\omega$ भी !!{enumerable} है}।
चूँकि हम जानते हैं कि $\Bin^\omega$ !!{enumerable} नहीं है, इसलिए
$\Pow{\PosInt}$ भी गणनीय नहीं हो सकता। इसे एक समस्या का दूसरी
समस्या में \emph{न्यूनीकरण} कहते हैं---यहाँ हम $\Bin^\omega$ को गिनने की
समस्या का न्यूनीकरण $\Pow{\PosInt}$ को गिनने की समस्या में करते हैं। दूसरी
समस्या का समाधान---$\Pow{\PosInt}$ की गणना---पहली समस्या का समाधान,
अर्थात $\Bin^\omega$ की गणना, दे देता।

हम किसी समुच्चय~$B$ को गिनने की समस्या का न्यूनीकरण किसी समुच्चय~$A$ को
गिनने की समस्या में कैसे करें? हम~$A$ की किसी गणना को~$B$ की गणना में बदलने
का ढंग देते हैं। इसका सबसे सरल उपाय कोई !!a{surjective} फलन
$f\colon A \to B$ परिभाषित करना है। यदि $x_1$, $x_2$, \dots{},~$A$ को
गिनते हैं, तो $f(x_1)$, $f(x_2)$, \dots{},~$B$ को गिनेंगे। हमारे मामले में
हम कोई आच्छादक फलन $f\colon \Pow{\PosInt} \to \Bin^\omega$ खोज रहे हैं।

\begin{prob}
दिखाइए कि यदि कोई !!{injective} फलन $g\colon B \to A$ है और
$B$~!!{nonenumerable} है, तो~$A$ भी अगणनीय है। यह दिखाकर ऐसा
कीजिए कि~$g$ की सहायता से~$A$ की गणना को~$B$ की गणना में कैसे बदला जा सकता
है।
\end{prob}

\begin{proof}[न्यूनीकरण द्वारा {\olref[nen]{thm:nonenum-pownat}} का प्रमाण]
मान लीजिए कि $\Pow{\PosInt}$ !!{enumerable} है और इसलिए उसकी कोई गणना
$Z_{1}$, $Z_{2}$, $Z_{3}$, \dots{} है।

फलन $f \colon \Pow{\PosInt} \to \Bin^\omega$ इस प्रकार परिभाषित करें:
$f(Z)$ वह अनुक्रम $s_{k}$ हो जिसके लिए $s_{k}(n) = 1$ यदि और केवल यदि
$n \in Z$, और अन्यथा $s_k(n) = 0$। इससे निश्चय ही एक फलन परिभाषित होता है,
क्योंकि जब भी $Z \subseteq \PosInt$ हो, तब कोई भी $n \in \PosInt$ या तो
$Z$ का !!a{element} है या नहीं है। उदाहरण के लिए, धन सम संख्याओं का समुच्चय
$2\PosInt = \{2,
4, 6, \dots\}$ अनुक्रम $010101\dots$ पर भेजा जाता है,
रिक्त समुच्चय $0000\dots$ पर और स्वयं $\PosInt$, $1111\dots$ पर भेजा जाता
है।

यह फलन !!{surjective} भी है: $0$ और $1$ का प्रत्येक अनुक्रम धन पूर्णांकों के
किसी समुच्चय के संगत है, अर्थात उस समुच्चय के जिसके अवयव वे पूर्णांक हैं जो
अनुक्रम में $1$ आने वाले स्थानों के संगत हैं। अधिक ठीक रूप में, मान लीजिए
$s \in \Bin^\omega$। समुच्चय $Z
\subseteq \PosInt$ इस प्रकार परिभाषित
करें:
\[
Z = \Setabs{n \in \PosInt}{s(n) = 1}
\]
तब $f(Z) = s$; इसे~$f$ की परिभाषा देखकर जाँचा जा सकता है।

अब सूची
\[
f(Z_1), f(Z_2), f(Z_3), \dots
\]
पर विचार कीजिए। चूँकि $f$ !!{surjective} है, इसलिए $\Bin^\omega$ का प्रत्येक
अवयव किसी तर्क के लिए~$f$ के मान के रूप में अवश्य आता है और इस कारण सूची में
भी अवश्य आता है। अतः यह सूची पूरे~$\Bin^\omega$ को गिनती है।

यदि $\Pow{\PosInt}$ !!{enumerable} होता, तो $\Bin^\omega$ भी
!!{enumerable} होता। पर $\Bin^\omega$ !!{nonenumerable} है
(\olref[nen]{thm:nonenum-bin-omega})। अतः $\Pow{\PosInt}$
!!{nonenumerable} है।
\end{proof}

\begin{explain}
न्यूनीकरण किस दिशा में होता है, इसे लेकर भ्रम होना आसान है। उदाहरण के लिए,
कोई !!a{surjective} फलन $g \colon \Bin^\omega \to B$ यह \emph{सिद्ध नहीं}
करता कि $B$ !!{nonenumerable} है। (उस फलन $g
\colon \Bin^\omega \to \Bin$
पर विचार कीजिए जिसे $g(s) = s(1)$ द्वारा परिभाषित किया गया है: यह फलन
$0$ और $1$ के किसी अनुक्रम को उसके पहले !!{element} पर भेजता है। यह
!!{surjective} है, क्योंकि कुछ अनुक्रम $0$ से और कुछ $1$ से आरंभ होते हैं।
लेकिन $\Bin$ परिमित है।) यह भी ध्यान रखें कि फलन~$f$ का !!{surjective} होना
आवश्यक है; अन्यथा तर्क पूरा नहीं होता: तब यह गारंटी नहीं होती कि $f(x_1)$,
$f(x_2)$, \dots{} में~$B$ के सभी !!{element} सम्मिलित हैं। उदाहरण के लिए,
\[
h(n) = \underbrace{000\dots0}_{\text{$n$ बार $0$}}
\]
एक फलन $h\colon \PosInt \to
\Bin^\omega$ परिभाषित करता है, लेकिन
$\PosInt$ !!{enumerable} है।
\end{explain}

\begin{prob}
न्यूनीकरण के तर्क द्वारा दिखाइए कि धन पूर्णांकों के युग्मों के सभी
\emph{समुच्चयों का समुच्चय} !!{nonenumerable} है।
\end{prob}

\begin{prob}\label{sfr:siz:red:prob:nat-nat}
  न्यूनीकरण के तर्क द्वारा दिखाइए कि समुच्चय~$X$, जिसमें सभी फलन
  $f\colon \Nat \to \Nat$ हैं, !!{nonenumerable} है (संकेत: $X$
  से~$\Bin^\omega$ पर कोई
  आच्छादक फलन दीजिए।)
\end{prob}

\begin{prob}
न्यूनीकरण के तर्क द्वारा दिखाइए कि $\Nat^\omega$, अर्थात प्राकृत संख्याओं के
सभी अनंत अनुक्रमों का समुच्चय, !!{nonenumerable} है।
\end{prob}

\begin{prob}
$P$, धन पूर्णांकों के समुच्चय से समुच्चय $\{0\}$ में जाने वाले फलनों का
समुच्चय हो, और $Q$, धन पूर्णांकों के समुच्चय से समुच्चय $\{0\}$ में जाने
वाले \emph{आंशिक} फलनों का समुच्चय हो। दिखाइए कि $P$~!!{enumerable} है और
$Q$~गणनीय नहीं है। (संकेत: $\Bin^\omega$ को गिनने की समस्या का
न्यूनीकरण~$Q$ को गिनने की समस्या में कीजिए।)
\end{prob}

\begin{prob}
$S$, धन पूर्णांकों के समुच्चय से समुच्चय \{0,1\} पर जाने वाले सभी
!!{surjective} फलनों का समुच्चय हो; अर्थात $S$ में सभी !!{surjective} फलन
$f\colon \PosInt \to \Bin$ हैं। दिखाइए कि $S$ !!{nonenumerable} है।
\end{prob}

\begin{prob}
दिखाइए कि सभी वास्तविक संख्याओं का समुच्चय~$\Real$ !!{nonenumerable} है।
\end{prob}

\end{document}
